\section{Background}
\label{section:background}
\vspace{-1ex}
% \subsection{Problem Definition}
% \textbf{Notations}\quad
\paragraph{Notations.} We represent each molecule as a graph $\gG =  (\gV, \gE)$, where $\gV$ is the set of vertices representing atoms and $\gE$ is the set of edges representing inter-atomic bonds. Each node $v$ in $\gV$ describes the chosen atomic features such as element type, atomic charge, and hybridization state. Each edge $e_{uv}$ in $\gE$ describes the corresponding chemical bond connecting $u$ and $v$, and is labeled with its bond type. Following \citet{GraphDG}, each molecular graph is expanded to incorporate auxiliary edges connecting all atoms within a 4\AA\ radius to enhance long-range interactions in message passing.
The spatial position of each atom in $\gV$ is represented by a 3D coordinate vector $\vr \in \sR^{3}$, such that the full molecule conformation is represented by the matrix $\mX \in \sR^{|\gV| \times 3}$. 

\vspace{-1ex}
\paragraph{Problem Definition.} \emph{Molecular conformation generation} (MCG) is a conditional generative process that aims to model the conditional distribution of 3D molecular conformations $\mX$, given the 2D molecule graph $\gG$, \ie, $p(\mX | \gG)$. While prior works have shown some success with learning 3D conformations starting with only the 2D information, many require complex and compute-intensive architectures \citep{DMCG, unimol}. Recently, \citet{torsional-diffusion} demonstrated good performance on the GEOM-DRUGS dataset by priming the method with easy-to-obtain 3D approximations via RDKit EKTDG \citep{rdkit}. \citet{torsional-diffusion} showed that RDKit is highly effective at generating conformations with correct inter-atomic distances and, as a result, can constrain the problem to a diffusion process only over torsion angles.

We thus formalize MCG as modeling the conditional distribution $p(\mX | \gR)$, where $\gR$ = $\langle \gV, \gE, \gX \rangle$ and $\gX$ is the RDKit generated atomic coordinates. This is functionally the same underlying distribution as $p(\mX | \gG)$, as we use RDKit as a building block to provide an approximation starting from only 2D information.
We will show that more robust conformers are generated by conditioning on approximations without imposing explicit angular and distance constraints. 

\vspace{-1ex}
\paragraph{Classical Methods for Conformer Generation.} A molecular conformer refers to the collection of 3D structures that are energetically favorable and correspond to local minima of the potential energy surface.
CREST~\citep{CREST} uses semi-empirical tight-binding density functional theory for energy calculations, which, while computationally less expensive than ab-initio quantum mechanical (QM) methods, still requires approximately 90 core hours per drug-like molecule~\citep{GEOM}. Though CREST was used to generate the ground truth GEOM dataset, it is too slow for downstream applications such as high-throughput virtual screening.

Cheminformatics methods, such as RDKit EKTDG, are commonly used to quickly generate approximate low-energy conformations of molecules. 
These methods are less accurate than QM methods due to the sparse coverage of the conformational space resulting from stochastic sampling. 
Additionally, force field optimizations are inherently less accurate than the above QM methods. RDKit EKTDG employs a genetic algorithm for Distance Geometry optimization that can be enhanced with a molecular mechanics force field optimization (MMFF).
\citet{rdkit-cluster} recently highlighted potential flaws in the current MCG benchmarks and demonstrated how EKTDG, in conjunction with a traditional sampling and clustering procedure, can outperform many recent deep learning methods.
To overcome the accuracy-efficiency trade-off in classical methods, we combine easy-to-obtain classical approximations with efficient data-driven strategies to improve the generalization of robust MCG.
\vspace{-1ex}
\paragraph{Deep Learning Methods for Conformer Generation.} Several probabilistic deep learning methods for MCG have been developed~\citep{generative_background}, such as variational autoencoders in CVGAE \citep{CVGAE} and ConfVAE \citep{ConfVAE}, normalizing flows in CGCF \citep{CGCF}, score-based generative models in ConfGF \citep{ConfGF} and DGSM \citep{DGSM}, and diffusion models in GeoDiff \citep{geodiff} and Torsional Diffusion \citep{torsional-diffusion}. GraphDG \citep{GraphDG} forgoes modeling coordinates and angles, relying solely on distance geometry. DMCG \citep{DMCG} and Uni-Mol \citep{unimol} present examples of effective large models, the first mimicking the architecture of AlphaFold \citep{alphafold} and the second using large-scale SE(3)-equivariant transformer pretraining. 

%Outside of MCG, there are also various deep learning methods that focus on generating valid and stable 3D molecules from scratch~\citep{simm2020symmetry, luo2022an, hoogeboom2022equivariant, xu2023geometric, qiang2023coarse}

\vspace{-1ex}
\paragraph{Molecular Coarse-graining.} Molecular coarse-graining refers to the simplification of a molecule representation by grouping the fine-grained (FG) atoms in the original structure into individual coarse-grained (CG) beads $\gB$ with a rule-based mapping.
Coarse-graining has been widely utilized in protein design \citep{CGProt-background, CGProtein} and molecular dynamics~\citep{cg-force-field}, and analogously fragment-level or subgraph-level generation has proven to be highly valuable in diverse 2D molecule design tasks \citep{FaST}.
Breaking down generative problems into smaller pieces is an approach that can be applied to several 3D molecule tasks. For instance, CGVAE \citep{CGVAE} learns a latent distribution to back map or restore FG coordinates from a fixed number of CG beads effectively.
We note that various different coarse-graining strategies exist~\citep{cg-survey, arts2023diffusion, CG-MD, cg_thermo}, and many require the ability to represent inputs with a non-fixed granularity. 
To handle this, \model uses a flexible variable-length CG representation that is compatible with all coarse-graining techniques.

\vspace{-1ex}
\paragraph{Autoregressive Molecule Generation.} Autoregressive models provide control over the generative process by enabling direct conditioning on prior information, allowing for a more precise and targeted generation of output.
Autoregressive generation has shown success in 2D molecule tasks using SMILE-based methods, as seen in MolMIM~\citep{reidenbach2023improving}, as well as graph-based atom-wise and subgraph-level techniques, as shown in GraphAF \citep{graphaf} and HierVAE \citep{HierVAE}.
Similarly, 3DLinker \citep{3dlinker} and SQUID \citep{squid} showcase the usefulness of 3D autoregressive molecule generation and their ability to leverage conditional information in both atom-wise and subgraph-level settings for 3D linkage and shape-conditioned generative tasks respectively. 
 We note that, unlike prior methods~\citep{squid}, \model does not require a predefined fragment vocabulary.
HERN~\citep{jin2022antibody} further demonstrates the power of hierarchical equivariant autoregressive methods in the task of computational 3D antibody design.

\vspace{-1ex}
\paragraph{SE(3)-Equivariance.} Let $\gX$ and $\gY$ be the input and output vector spaces, respectively, which possess a set of transformations $G$: $G\times\gX\rightarrow\gX$ and $G\times\gY\rightarrow\gY$. The function $\phi: \gX\rightarrow\gY$ is called equivariant with respect to $G$ if when we apply any transformation to the input, the output also changes via the same transformation or under a certain predictable behavior, \ie,
\begin{definition}
The function $\phi: \gX\mapsto\gY$ is $G$-equivariant if it commutes with any transformation in $G$,
\begin{eqnarray}
\label{eq:equ}
\phi (\rho_{\gX}(g)x) = \rho_{\gY}(g)\phi(x), \forall g\in G,
\end{eqnarray}
where $\rho_{\gX}$ and $\rho_{\gY}$ are the group representations in the input and output space, respectively. Specifically, $\phi$ is called invariant if $\rho_\gY$ is the identity.
\end{definition}
By enforcing SE(3)-equivariance in our probabilistic model,  $p(\mX | \gR)$ remains unchanged for any rototranslation of the approximate conformer $\gR$.
\modelns's architecture is inspired by recent equivariant graph neural network architectures, such as EGNN \citep{EGNN} and PaiNN \citep{painn}, as well as Vector Neuron multi-layer perceptron (VN-MLP) \citep{vn-mlp}.
\vspace{-2ex}